%============================================================
%  ch09_TB_manifold.tex
%  Chapter 9: The Tree-Based Manifold
%
%  Tensor Formats in Banach Spaces
%  Antonio Falcó, CEU Universities
%
%  Based on: FHN2023 (Secs. 4-5)
%  Estimated length: ~40 pp.
%============================================================

\chapter{The Tree-Based Manifold}
\label{chap:TB_manifold}

This chapter establishes the geometric counterpart of
Chapter~\ref{chap:tucker_manifold} for general tree-based tensor
formats. The central result is that $\FTv$ — when it is a
\emph{proper} set of tree-based tensors (Definition~\ref{def:proper_TBT})
— carries a natural $\mathcal{C}^\infty$-Banach manifold structure
that is compatible with the Tucker manifold structures at each
level of the tree, and is an immersed submanifold of the ambient
Banach tensor space $\VDnorm$.

The key idea, due to~\cite{FHN2023}, is that
$\FTv = \bigcap_{k=1}^{\dep{\TD}} \Tucker{\fr_k}{\VD, \PkT{k}{\TD}}$
(Theorem~\ref{thm:intersection_representation}) is an intersection of
Tucker manifolds, and its chart system is obtained by taking the
intersection of the Tucker chart systems at each level. The resulting
charts are natural and compatible with those of each Tucker factor.

Throughout, $(V_\alpha, \|\cdot\|_\alpha)$ are normed spaces,
$\|\cdot\|_D$ satisfies \condINJ\ at each level of $\TD$, and
$\fr \in \ADTD$.

%------------------------------------------------------------
\section{Proper sets of tree-based tensors}
\label{sec:TB_manifold:proper}
%------------------------------------------------------------

Not every $\FTv$ is a genuine intersection of Tucker manifolds — in
degenerate cases it may coincide with a single Tucker manifold.
The following definition isolates the non-degenerate case.

\begin{definition}[Proper tree-based set]
\label{def:proper_TBT}
The set $\FTv$ is called a \emph{proper set of tree-based tensors}
with rank $\fr$ if
\begin{equation*}
    \FTv \neq \Tucker{\fr_k}{\VD, \PkT{k}{\TD}}
    \quad \text{for all } 1 \leq k \leq \dep{\TD}.
\end{equation*}
\end{definition}

\begin{example}[Degenerate case]
\label{ex:degenerate}
For $D = \{1,\ldots,6\}$ and the tree $\TD$ with
$\Sons{D} = \{\{1\}, \{2,3,4,5,6\}\}$ (depth~2):
\begin{align*}
    \mathcal{P}_1(\TD) &= \bigl\{\{1\}, \{2,3,4,5,6\}\bigr\}, \\
    \mathcal{P}_2(\TD) &= \Leaves.
\end{align*}
Since $\dim \Umin{\{1\}}{\bv} = \dim \Umin{\{2,3,4,5,6\}}{\bv}$ for all
$\bv \in \VD$ (Theorem~\ref{thm:min_existence}), the constraint at level~1
is automatically satisfied whenever the leaf-level constraints hold.
Hence $\FTv = \Tucker{\fr_2}{\VD, \Leaves}$ and $\FTv$ is \emph{not}
a proper set. The geometry reduces to the Tucker manifold.
\end{example}

In all examples arising from balanced binary trees (HT format),
Tensor Train, and most practical trees with $\dep{\TD} \geq 2$,
the set $\FTv$ is proper.

%------------------------------------------------------------
\section{The Laplacian-like map and the intersection of model spaces}
\label{sec:TB_manifold:laplacian}
%------------------------------------------------------------

For each level $k$ of the tree, the Tucker manifold
$\Tucker{\fr_k}{\VD, \PkT{k}{\TD}}$ has model space
\begin{equation*}
    \Emodel{k} \coloneqq
    \Bigl(\bigoplus_{\alpha \in \PkT{k}{\TD}}
    \cL(\Umin{\alpha}{\bv}, \Wmin{\alpha}{\bv})
    \otimes \spann\{\mathbf{id}_{[\alpha]}\}\Bigr)
    \times \RR^{\times_{\alpha \in \PkT{k}{\TD}} r_\alpha},
\end{equation*}
where $\mathbf{id}_{[\alpha]} = \bigotimes_{\beta \in \PkT{k}{\TD}
\setminus \{\alpha\}} \Id{\Vbigbeta}$ is the identity on all factors
except $\alpha$.

The \emph{Laplacian-like map} $\Delta_k$ embeds the product of
$\cL(\Umin{\alpha}{\bv}, \Wmin{\alpha}{\bv})$ spaces into $\cL(\VDnorm)$:
\begin{equation}
\label{eq:laplacian}
    \Delta_k : \prod_{\alpha \in \PkT{k}{\TD}}
    \cL(\Umin{\alpha}{\bv}, \Wmin{\alpha}{\bv})
    \longrightarrow \cL(\VDnorm),
    \quad
    (L_\alpha)_\alpha \mapsto \sum_{\alpha \in \PkT{k}{\TD}}
    L_\alpha \otimes \mathbf{id}_{[\alpha]}.
\end{equation}

\begin{proposition}[{\cite[Proposition~4.5]{FHN2023}}]
\label{prop:laplacian_exp}
For $(L_\alpha)_{\alpha \in \PkT{k}{\TD}} \in \prod_\alpha
\cL(\Umin{\alpha}{\bv}, \Wmin{\alpha}{\bv})$:
\begin{equation}
\label{eq:exp_laplacian}
    \exp\bigl(\Delta_k((L_\alpha)_\alpha)\bigr)
    = \bigotimes_{\alpha \in \PkT{k}{\TD}} \exp(L_\alpha).
\end{equation}
\end{proposition}

\begin{proof}
Put $L = \sum_\alpha L_\alpha \otimes \mathbf{id}_{[\alpha]}$.
For $\alpha \neq \beta$:
$(L_\alpha \otimes \mathbf{id}_{[\alpha]}) \circ
(L_\beta \otimes \mathbf{id}_{[\beta]}) =
L_\alpha \otimes L_\beta \otimes \mathbf{id}_{[{\alpha,\beta}]}$,
and by the nilpotency of $L_\alpha$ (Proposition~\ref{prop:nilpotent_algebra}),
$\exp(L_\alpha \otimes \mathbf{id}_{[\alpha]}) = \exp(L_\alpha) \otimes
\mathbf{id}_{[\alpha]}$. Since the summands commute,
$\exp(L) = \bigodot_\alpha \exp(L_\alpha \otimes \mathbf{id}_{[\alpha]})
= \bigotimes_\alpha \exp(L_\alpha)$.
\end{proof}

The critical consequence is that the local representation
$\bw = \exp(L)(\bu)$ for $L \in \Emodel{k}$ is
\emph{independent of $k$}: for any $k$ and any $L \in \Emodel{k}$
with the same underlying tensor $L = \Delta_k((L_\alpha)_\alpha)$,
the formula $\bw = (\bigotimes_\alpha \exp(L_\alpha))(\bu)$ gives
the same $\bw$.

%------------------------------------------------------------
\section{Charts for the tree-based manifold}
\label{sec:TB_manifold:charts}
%------------------------------------------------------------

\subsection*{Intersection of chart domains}

For $\bv \in \FTv$, define the \emph{tree-based chart domain}:
\begin{equation}
\label{eq:TBT_chart_domain}
    \mathcal{W}(\bv) \coloneqq
    \bigcap_{k=1}^{\dep{\TD}} \chart^{(k)}(\bv),
\end{equation}
where $\chart^{(k)}(\bv)$ is the chart domain of
$\Tucker{\fr_k}{\VD, \PkT{k}{\TD}}$ at $\bv$
(Section~\ref{sec:tucker_manifold:charts}).

\subsection*{Model space}

Define the \emph{tree-based model space}:
\begin{equation}
\label{eq:TBT_model}
    \mathbf{E}(\bv) \coloneqq
    \bigcap_{k=1}^{\dep{\TD}} \Emodel{k}
    \subset \cL(\VDnorm),
\end{equation}
and the \emph{tree-based core set}:
\begin{equation*}
    \mathcal{O}(\bv) \coloneqq
    \bigcap_{k=1}^{\dep{\TD}}
    \Tucker{\fr_k}\Bigl({}_a\bigotimes_{\alpha \in \PkT{k}{\TD}}
    \Umin{\alpha}{\bv}\Bigr).
\end{equation*}

\begin{lemma}[{\cite[Lemma~4.6]{FHN2023}}]
\label{lem:TBT_chart_open}
Under condition \condINJ\ at each level, $\mathbf{E}(\bv) \times
\mathcal{O}(\bv)$ is an open subset of the Banach space
$\mathbf{E}(\bv) \times {}_a\bigotimes_{\alpha \in \PkT{1}{\TD}}
\Umin{\alpha}{\bv}$.
\end{lemma}

\begin{proof}[Sketch]
Each $\mathcal{O}_k \coloneqq \Tucker{\fr_k}({}_a\bigotimes_\alpha
\Umin{\alpha}{\bv})$ is open in its finite-dimensional ambient space
(it is defined by rank conditions, cf.\ Remark~\ref{rem:Rstar_open}).
The intersection $\mathcal{O}(\bv) = \bigcap_k \mathcal{O}_k$ is
taken in ${}_a\bigotimes_{\alpha \in \PkT{1}{\TD}} \Umin{\alpha}{\bv}$
(since all factors have the same underlying space by the tree structure),
and is an open subset. Each $\Emodel{k}$ is a closed subspace of
$\cL(\VDnorm)$ (Lemma~\ref{lem:two_faces}), so $\mathbf{E}(\bv)$ is
a Banach space and $\mathbf{E}(\bv) \times \mathcal{O}(\bv)$ is open.
\end{proof}

\subsection*{Tree-based chart map}

Since $L \in \mathbf{E}(\bv)$ satisfies $L \in \Emodel{k}$ for all
$k$, the formula $\bw = \exp(L)(\bu)$ is well-defined and independent
of the level index $k$ by Proposition~\ref{prop:laplacian_exp}.
This allows us to define:

\begin{definition}[Tree-based chart]
\label{def:TBT_chart}
For $\bv \in \FTv$, the \emph{tree-based chart map}
\begin{equation}
\label{eq:TBT_chart_map}
    \boldsymbol{\xi}_\bv : \mathcal{W}(\bv)
    \longrightarrow \mathbf{E}(\bv) \times \mathcal{O}(\bv)
\end{equation}
is defined by $\boldsymbol{\xi}_\bv(\bw) = (L, \bu)$ where
$\bw = \exp(L)(\bu)$, with $L \in \mathbf{E}(\bv)$ and
$\bu \in \mathcal{O}(\bv)$ uniquely determined by any Tucker chart
$\chartmap{\bv}^{(k)}$.
\end{definition}

%------------------------------------------------------------
\section{The $\mathcal{C}^\infty$-Banach manifold structure}
\label{sec:TB_manifold:cinfty}
%------------------------------------------------------------

\begin{theorem}[$\mathcal{C}^\infty$ tree-based manifold,
{\cite[Theorem~4.7]{FHN2023}}]
\label{thm:TBT_Banach_Manifold}
Let $\TD$ be a dimension partition tree over $D$ with $\dep{\TD} \geq 2$,
and let $\fr \in \ADTD$ be such that $\FTv$ is a proper set of
tree-based tensors. Assume \condINJ\ holds at each level of $\TD$.
Then:
\begin{enumerate}
\item[\textup{(i)}] The collection
    $\{(\mathcal{W}(\bv), \boldsymbol{\xi}_\bv)\}_{\bv \in \FTv}$
    is a $\mathcal{C}^\infty$-atlas on $\FTv$, making it a
    $\mathcal{C}^\infty$-Banach manifold modelled on
    $\mathbf{E}(\bv) \times {}_a\bigotimes_{\alpha \in \PkT{1}{\TD}}
    \Umin{\alpha}{\bv}$.

\item[\textup{(ii)}] The atlas is compatible with the Tucker manifold
    atlas at each level: the inclusion
    \begin{equation*}
        \mathfrak{i}_k : \FTv \hookrightarrow
        \Tucker{\fr_k}{\VD, \PkT{k}{\TD}}
    \end{equation*}
    is a $\mathcal{C}^\infty$-immersion for each $1 \leq k \leq \dep{\TD}$.
\end{enumerate}
\end{theorem}

\begin{proof}[Sketch]
The atlas conditions follow from Lemma~\ref{lem:TBT_chart_open}
and the fact that the transition maps of the tree-based atlas are
restrictions of the transition maps of the Tucker atlases (which are
$\mathcal{C}^\infty$ by Lemma~\ref{lem:premanifold}). Compatibility
with the Tucker atlases is a consequence of the identity
$\boldsymbol{\xi}_\bv = \chartmap{\bv}^{(k)}|_{\mathcal{W}(\bv)}$
for each $k$. See~\cite{FHN2023} for the complete proof.
\end{proof}

%------------------------------------------------------------
\section{Tangent space of the tree-based manifold}
\label{sec:TB_manifold:tangent}
%------------------------------------------------------------

\begin{proposition}[Tangent space]
\label{prop:TBT_tangent}
Under the assumptions of Theorem~\ref{thm:TBT_Banach_Manifold},
for $\bv \in \FTv$:
\begin{equation}
\label{eq:TBT_tangent}
    T_\bv \mathfrak{i}\bigl(\mathbb{T}_\bv(\FTv)\bigr)
    = \bigcap_{k=1}^{\dep{\TD}}
    T_\bv \mathfrak{i}_k
    \bigl(\mathbb{T}_\bv(\Tucker{\fr_k}{\VD, \PkT{k}{\TD}})\bigr)
    = \bigcap_{k=1}^{\dep{\TD}} \ZD[\PkT{k}{\TD}]{\bv},
\end{equation}
where $\ZD[\mathcal{P}]{\bv}$ is the tangent subspace
\eqref{eq:ZD} for the partition $\mathcal{P}$.
\end{proposition}

\begin{proof}
The model space of $\FTv$ is $\mathbf{E}(\bv) \times \mathcal{O}_1$,
where $\mathbf{E}(\bv) = \bigcap_k \Emodel{k}$. Since $T_\bv \mathfrak{i}_k$
is injective (Proposition~\ref{prop:tangent_space}), the tangent map
$T_\bv \mathfrak{i}$ is also injective and its image is
$\bigcap_k T_\bv \mathfrak{i}_k(\mathbb{T}_\bv \Tucker{\fr_k}{\VD, \PkT{k}})$.
The last equality follows from the explicit formula
\eqref{eq:tangent_formula} for each level.
\end{proof}

%------------------------------------------------------------
\section{Immersion in the ambient Banach space}
\label{sec:TB_manifold:immersion}
%------------------------------------------------------------

\begin{theorem}[Immersion of the tree-based manifold,
{\cite[Corollary~4.9]{FHN2023}}]
\label{thm:TBT_immersion}
Under the assumptions of Theorem~\ref{thm:TBT_Banach_Manifold}:
\begin{equation}
\label{eq:TBT_tangent_split}
    T_\bv \mathfrak{i}\bigl(\mathbb{T}_\bv(\FTv)\bigr)
    \in \GrassB{\VDnorm}
    \quad \text{for all } \bv \in \FTv.
\end{equation}
Consequently, $\FTv$ is an immersed submanifold of $\VDnorm$.
\end{theorem}

\begin{proof}
From Theorem~\ref{thm:tucker_immersion} applied at each level,
$\ZD[\PkT{k}]{\bv} \in \GrassB{\VDnorm}$ for $1 \leq k \leq \dep{\TD}$.
By Proposition~\ref{prop:TBT_tangent}, the tangent space of $\FTv$
is $\bigcap_k \ZD[\PkT{k}]{\bv}$. Applying
Lemma~\ref{lem:direct_sum_G} iteratively (intersection of split
subspaces is split), we obtain
$\bigcap_k \ZD[\PkT{k}]{\bv} \in \GrassB{\VDnorm}$.
\end{proof}

\begin{remark}[Embedded vs.\ immersed]
\label{rem:embedded_vs_immersed}
Theorem~\ref{thm:TBT_immersion} shows that $\FTv$ is an
\emph{immersed} submanifold of $\VDnorm$. Moreover,
Theorem~\ref{thm:TBT_Banach_Manifold}(ii) shows it is an
\emph{embedded} submanifold of each Tucker manifold
$\Tucker{\fr_k}{\VD, \PkT{k}{\TD}}$. Whether $\FTv$ is embedded
in $\VDnorm$ (i.e., the manifold topology coincides with the
subspace topology) is a more delicate question related to the
global topology discussed in Section~\ref{sec:TB_manifold:topology}.
\end{remark}

%------------------------------------------------------------
\section{Comparison with alternative descriptions}
\label{sec:TB_manifold:comparison}
%------------------------------------------------------------

The geometric description of $\FTv$ in~\cite{FHN2023} differs from
earlier approaches in several respects.

\paragraph{Uschmajew--Vandereycken \cite{UschmajewVandereycken2013}.}
For the HT format in finite-dimensional spaces, they construct
$\FTv$ as a quotient manifold
$P \backslash \bigl(\GL{r_\alpha} \times \cdots\bigr)$
for a gauge group $P$. This gives a finite-dimensional manifold
and is well-suited for numerical algorithms (Riemannian gradient
methods). However, it does not extend to infinite-dimensional
spaces and does not make the compatibility with Tucker manifolds
at each level explicit.

\paragraph{Holtz--Rohwedder--Schneider \cite{HoltzRohwedderSchneider2012}.}
For the TT format in finite dimensions, a similar quotient
construction is used. The connection to the Tucker manifold at
each level is not made explicit.

\paragraph{Falcó--Hackbusch--Nouy (2015 preprint \cite{FHN2015}).}
An earlier version of the atlas in~\cite{FHN2023} was proposed
using a different chart system. The description in~\cite{FHN2023}
is more natural because the chart maps are directly inherited from
the Tucker manifold charts via the Laplacian map (Proposition~\ref{prop:laplacian_exp}).

\paragraph{Key novelties of \cite{FHN2023}.}
\begin{enumerate}
\item The atlas $\{(\mathcal{W}(\bv), \boldsymbol{\xi}_\bv)\}$
    is \emph{natural}: it is the restriction of the Tucker atlases
    to the intersection $\FTv$, without any additional construction.
\item The \emph{compatibility} with Tucker manifold structures at
    each level is built into the definition, making the geometric
    relationships transparent.
\item The description holds in \emph{infinite-dimensional Banach
    spaces} under the mild condition \condINJ.
\item The tangent space formula \eqref{eq:TBT_tangent} is explicit
    and computable.
\end{enumerate}

%------------------------------------------------------------
\section{Topology and connected components}
\label{sec:TB_manifold:topology}
%------------------------------------------------------------

\begin{proposition}
\label{prop:TBT_connected}
For the Tucker manifold, $\Tucker{\fr}{\VD}$ is connected for any
$\fr \in \AD{\VD}$ with $\fr \neq \mathbf{0}$. This follows because
$\prod_\alpha \Grass{r_\alpha}{V_\alpha}$ is connected (each
$\Grass{r}{V}$ is connected as a Banach manifold) and the fibre
$\RR^{\times r_\alpha}_*$ is connected.
\end{proposition}

For general tree-based formats, the connectedness of $\FTv$ is an
open problem in general. The following is known:

\begin{proposition}
\label{prop:TBT_connected_cases}
\begin{enumerate}
\item[\textup{(i)}] For the TT format and binary HT formats in
    \emph{finite dimensions}, $\FTv$ is connected
    (see~\cite{UschmajewVandereycken2013, HoltzRohwedderSchneider2012}).
\item[\textup{(ii)}] For general $\TD$ and infinite-dimensional
    spaces, connectedness of $\FTv$ is an open question.
\end{enumerate}
\end{proposition}

\begin{remark}[Stratification]
\label{rem:stratification}
The decomposition
$\VD = \bigsqcup_{\fr \in \ADTD} \FT{\fr}{\VD}{\TD}$
(Theorem~\ref{thm:TBT_decomposition}) is a stratification of $\VD$
by tree-based rank. Each stratum $\FTv$ is an immersed manifold,
and $\FTleq{\fr}{\VD}{\TD} = \overline{\FTv}^{\text{weak}}$
is the weak closure. The boundary strata
$\partial \FTleq{\fr}{\VD}{\TD} = \FTleq{\fr}{\VD}{\TD} \setminus \FTv$
consist of tensors with strictly smaller rank.
This stratification is the tensor-format analogue of the rank
stratification of matrices.
\end{remark}

%------------------------------------------------------------
\section{Historical notes}
\label{sec:TB_manifold:history}
%------------------------------------------------------------

The geometric structure of $\FTv$ as a Banach manifold was established
in~\cite{FHN2023}. The fundamental observation — that the tree-based
manifold is an intersection of Tucker manifolds and that its atlas can
be obtained by intersecting the Tucker atlases — was made possible by
the reformulation of the Tucker manifold charts using the Laplacian-like
map $\Delta_k$ (Proposition~\ref{prop:laplacian_exp}).

The compatibility condition between the tree-based atlas and the Tucker
atlases at each level (Theorem~\ref{thm:TBT_Banach_Manifold}(ii)) is
the key property that distinguishes the description of~\cite{FHN2023}
from earlier work: it allows one to apply the Tucker manifold geometry
(tangent spaces, projections, Dirac--Frenkel principle) at each level
of the tree independently, and then combine the results.

The Dirac--Frenkel variational principle for tree-based formats
(Chapter~\ref{chap:DF}) exploits this compatibility in an essential
way: the reduced ODE on $\FTv$ is obtained by projecting the vector
field onto the tangent space $\bigcap_k \ZD[\PkT{k}]{\bv}$, and the
existence of this projection (Corollary~\ref{cor:best_proj_tangent}
applied at each level) follows from the fact that each $\ZD[\PkT{k}]{\bv}$
is a split subspace of $\VDnorm$.

